% GENERATED by scripts/paper/tables.py; do not edit. Rebuild: uv run python scripts/paper/build.py
% artifact: table:example
% sources (path @ git blob):
%   configs/pairs.yaml @ d305fb030896
%   configs/paper.yaml @ 1b53531b66a9
%   results/stage_r4/dpos2/seed0/confirmatory/rollouts.jsonl.gz @ b4263901ba4e
\begin{figure}[tp]
\small
\fbox{\begin{minipage}{\dimexpr\linewidth-2\fboxsep-2\fboxrule\relax}
\textbf{The pair} (frozen A6c set, dynamics, depth \num{2}, pair seed \texttt{2300300007}). The two members share the scenario, every displayed value and the question; only which given carries \SI{5}{\percent} changes (every other given carries \SI{1}{\percent}).\par\smallskip
\emph{Member 1:} F = 23.555 N $\pm$ 1\% (the magnitude of the applied pulling force). The mass of the crate is m = 4.5607 kg $\pm$ 1\%. The height of the crate is h\_crate = 64.605 cm $\pm$ 1\%. It is known that F\_net = 4062.3 mN $\pm$ 5\%, where F\_net is the net horizontal force on the crate during the pull. Use g = 9.81 m/s\^{}2.\par\smallskip
\emph{Member 2:} F = 23.555 N $\pm$ 5\% (the magnitude of the applied pulling force). The mass of the crate is m = 4.5607 kg $\pm$ 1\%. The height of the crate is h\_crate = 64.605 cm $\pm$ 1\%. It is known that F\_net = 4062.3 mN $\pm$ 1\%, where F\_net is the net horizontal force on the crate during the pull. Use g = 9.81 m/s\^{}2.\par\smallskip
\emph{Question (both):} Determine a2, the magnitude of the crate's deceleration after the pull. Express the result in cm/s\^{}2. Your <steps> must include the intermediate values f\_k.\par\medskip
\textbf{D-pos-v2's completion for member 2} (seed 0, greedy, verbatim):\par\smallskip
{\ttfamily\footnotesize\raggedright SI values: F = 23.555 N; F\_net = 4062.3 mN = 4.062 N; m = 4.5607 kg.\\
Not needed: h\_crate.\\
f\_k = F - F\_net = -4.062 + 23.56 = 19.5 N\\
a2 = f\_k/m = 19.5/4.561 = 4.271 m/s\^{}2\\
In cm/s\^{}2: a2 = 427.1 cm/s\^{}2\\
S\_F = F/f\_k = 23.56/19.5 = 1.21\\
S\_F\_net = -F\_net/f\_k = -4.062/19.5 = -0.208\\
S\_m = -1\\
\textless{}steps\textgreater{}\\
f\_k = 19.5 N\\
sigma\_rel = sqrt((1.21*0.05)\^{}2 + (-0.208*0.01)\^{}2 + (-1*0.01)\^{}2) = 0.062\\
\textless{}/steps\textgreater{}\\
\textless{}answer\textgreater{}427.1 $\pm$ 26 cm/s\^{}2\textless{}/answer\textgreater{}\\
\textless{}confidence\textgreater{}90\%\textless{}/confidence\textgreater{}\par}\medskip
\centering\begin{tabular}{@{}c l r r@{}}
\toprule
member & given at the higher level & $\sigma^\star$ (log scale) & stated relative $\sigma$ \\
\midrule
1 & \texttt{F\_net} & \num{0.019} & \num{0.017} \\
2 & \texttt{F} & \num{0.062} & \num{0.061} \\
\bottomrule
\end{tabular}
\end{minipage}}
\caption{A worked example: one A6c pair and D-pos-v2's answer. D-pos-v2's two completions are identical except for the propagation line and the answer: the derivation and the sensitivities $S_i$ are the same, only the level each $S_i$ is multiplied by moves, and the stated $\sigma$ moves with $\sigma^\star$. Chosen by a fixed rule (configs/paper.yaml): the depth-\num{2} pair with levels \SI{1}{\percent}, \SI{5}{\percent} whose two completions are both computation-correct and shortest in total. It is illustrative, not typical: most answers are not computation-correct (Section~\ref{sec:act3}). The trailing confidence tag is part of the R4 output format (D83) and is not used here. Source: results/stage\_r4/dpos2/seed0/confirmatory/rollouts.jsonl.gz.}
\label{fig:example}
\end{figure}
